\documentclass[../../main.tex]{subfiles}

\begin{document}

\section{局部有限性}\label{section:ch06-locally-finite}

在本节中, 我们证明局部有限族的一些性质以及关于可度量化空间的一个关键性引理.

\begin{definition}\label{definition:ch06-locally-finite}
设 $X$ 是一个拓扑空间. $X$ 的一个子集族 $\mathcal{A}$ 称为局部有限的(locally finite), 若 $X$ 中的每一点都存在一个邻域只与 $\mathcal{A}$ 的有限多个成员相交.
\end{definition}

\begin{examples}
可以验证, 在拓扑空间 $\mathbb{R}$ 中, 区间族
\begin{align*}
\mathcal{A}=\{(n,n+2)\mid n\in\mathbb{Z}\}
\end{align*}
是局部有限的. 另一方面, 族
\begin{align*}
\mathcal{B}=\{(0,1/n)\mid n\in\mathbb{Z}_+\}
\end{align*}
在 $(0,1)$ 中是局部有限的, 但在 $\mathbb{R}$ 中却不是. 同样, 族
\begin{align*}
\mathcal{C}=\{(1/(n+1),1/n)\mid n\in\mathbb{Z}_+\}
\end{align*}
也是如此.
\end{examples}

\begin{lemma}\label{lemma:ch06-locally-finite-properties}
设 $\mathcal{A}$ 是 $X$ 的一个局部有限的子集族. 则
\begin{enumerate}[(a)]
\item $\mathcal{A}$ 的任意子族是局部有限的.
\item $\mathcal{A}$ 中元素的闭包构成的族 $\mathcal{B}=\{\overline{A}\}_{A\in\mathcal{A}}$ 是局部有限的.
\item $\overline{\bigcup_{A\in\mathcal{A}} A}=\bigcup_{A\in\mathcal{A}}\overline{A}$.
\end{enumerate}
\end{lemma}

\begin{proof}
(a) 是显然的. 以下证明(b). 注意, 与 $\overline{A}$ 相交的任何开集 $U$ 也必然与 $A$ 相交. 因此, 如果 $U$ 是 $x$ 的一个邻域, 它仅与 $\mathcal{A}$ 中有限多个成员相交, 那么 $U$ 最多只能与族 $\mathcal{B}$ 中同样个数的成员相交(可能与 $\mathcal{B}$ 中更少的成员相交, 因为即使 $A_1$ 与 $A_2$ 不等, $\overline{A_1}$ 与 $\overline{A_2}$ 相等也是可能的).

为了证明(c), 令 $Y$ 表示 $\mathcal{A}$ 的所有成员的并
\begin{align*}
\bigcup_{A\in\mathcal{A}}A=Y.
\end{align*}
易见, $\bigcup\overline{A}\subset\overline{Y}$. 在局部有限的条件下, 我们证明相反的包含关系也成立. 设 $x\in\overline{Y}$, $U$ 是 $x$ 的一个邻域, 它只与 $\mathcal{A}$ 中有限多个成员 $A_1,\cdots,A_k$ 相交. 我们将证明 $x$ 必属于 $\overline{A_1},\cdots,\overline{A_k}$ 中的某一个, 因而属于 $\bigcup\overline{A}$. 否则, 集合 $U-\overline{A_1}-\cdots-\overline{A_k}$ 就是 $x$ 的一个邻域, 它不与 $\mathcal{A}$ 中任何成员相交, 从而不与 $Y$ 相交, 这与假设 $x\in\overline{Y}$ 矛盾.
\end{proof}

对于 $X$ 的加标子集族也有与局部有限类似的概念. 加标族 $\{A_\alpha\}_{\alpha\in J}$ 称为局部有限的加标族(locally finite indexed family), 如果 $X$ 中的每一个元素 $x$ 都存在一个邻域仅和有限多个 $A_\alpha$ 相交. 局部有限性的这两种叙述之间有什么关系呢? 容易看到 $\{A_\alpha\}_{\alpha\in J}$ 为局部有限的加标族当且仅当作为集族它是局部有限的, 并且 $X$ 的每一个非空子集 $A$ 最多只对有限多个 $\alpha$ 有 $A$ 与 $A_\alpha$ 相等.

局部有限加标族只会在第 41 节中处理单位分拆时用到, 到时候我们再来关注局部有限加标族.

\begin{definition}\label{definition:ch06-countably-locally-finite}
$X$ 的子集族 $\mathcal{B}$ 称为可数局部有限的(countably locally finite), 若 $\mathcal{B}$ 可以表示成可数个局部有限族 $\mathcal{B}_n$ 的并.
\end{definition}

这一概念也叫做``$\sigma$ 局部有限'', 其中 $\sigma$ 来自测度论, 表示``可数并''. 注意, 可数族和局部有限族都是可数局部有限的.

\begin{definition}\label{definition:ch06-refinement}
设 $\mathcal{A}$ 是 $X$ 的子集族, $X$ 的子集族 $\mathcal{B}$ 称为 $\mathcal{A}$ 的加细(refinement)(或称为加细 $\mathcal{A}$), 若对于 $\mathcal{B}$ 中的每一个元素 $B$, 存在 $\mathcal{A}$ 的一个元素 $A$ 包含 $B$. 若 $\mathcal{B}$ 中的每一个元素都是开集, 则称 $\mathcal{B}$ 是 $\mathcal{A}$ 的开加细(open refinement). 若 $\mathcal{B}$ 中的每一个元素都是闭集, 则称其为 $\mathcal{A}$ 的闭加细(closed refinement).
\end{definition}

\begin{lemma}\label{lemma:ch06-metrizable-countable-locally-finite-refinement}
设 $X$ 是一个可度量化空间. 若 $\mathcal{A}$ 是 $X$ 的一个开覆盖, 则存在 $X$ 的可数局部有限的开覆盖 $\mathcal{E}$ 加细 $\mathcal{A}$.
\end{lemma}

\par
\begin{figure}[H]
\centering
\includegraphics[width=0.45\textwidth]{fig39-1.png}
\caption{}\label{figure:ch06-fig39-1}
\end{figure}

\par
\begin{figure}[H]
\centering
\includegraphics[width=0.45\textwidth]{fig39-2.png}
\caption{}\label{figure:ch06-fig39-2}
\end{figure}

\begin{proof}
为了证明这个引理, 我们要用到良序定理. 对于集族 $\mathcal{A}$ 选取一个良序 $<$, 用字母 $U$, $V$, $W$, $\cdots$ 表示 $\mathcal{A}$ 的元素.

选定 $X$ 的一个度量. 设 $n$ 是一个暂时固定的正整数. 给定 $\mathcal{A}$ 中的一个元素 $U$, 我们定义 $S_n(U)$ 为
\begin{align*}
S_n(U)=\{x\mid B(x,1/n)\subset U\}.
\end{align*}
它是 $U$ 的一个子集, 通过``缩紧''$U$ $1/n$ 的距离得到. ($S_n(U)$ 是一个闭集, 但这一点对于我们并不重要.) 现在, 我们用 $\mathcal{A}$ 的良序 $<$ 来讨论一个更小的集合. 对于 $\mathcal{A}$ 中的每一个 $U$, 定义
\begin{align*}
T_n(U)=S_n(U)-\bigcup_{V<U}V.
\end{align*}
如\cref{figure:ch06-fig39-1}所示, 其中 $\mathcal{A}$ 由三个集合 $U<V<W$ 组成. 这个图形正好反映出所得到的这些集合是无交的.


事实上, 可以断言, 它们至少被分开了 $1/n$ 的距离. 也就是说, 如果 $V$ 和 $W$ 是 $\mathcal{A}$ 中两个不同的元素, 那么对于任意 $x\in T_n(V)$ 和 $y\in T_n(W)$, 有 $d(x,y)\geqslant 1/n$.

为了证明这一事实, 设已选定 $V<W$. 由 $x\in T_n(V)$, 可得 $x\in S_n(V)$, 从而 $B(x,1/n)\subset V$. 另一方面, 因为 $V<W$, 并且 $y\in T_n(W)$, 根据 $T_n(W)$ 的定义可见 $y\notin V$, 于是 $d(x,y)\geqslant 1/n$.

然而集合 $T_n(U)$ 仍然不是我们想要的, 因为它不是开集(实际上它是闭集), 所以还需将其放大一点, 使之成为一个开集 $E_n(U)$. 特别地, 令 $E_n(U)$ 为 $T_n(U)$ 的 $\tfrac{1}{3n}$-邻域, 也就是说 $E_n(U)$ 是所有开球 $B\bigl(x,\tfrac{1}{3n}\bigr)$(其中 $x\in T_n(U)$)的并.

在 $U<V<W$ 的假设下, 如\cref{figure:ch06-fig39-2}所示, 我们所得到的集合是无交的. 事实上, 若 $V$ 和 $W$ 是 $\mathcal{A}$ 中不同的两个元素, 那么对于任意 $x\in E_n(V)$ 和 $y\in E_n(W)$, 一定有 $d(x,y)\geqslant 1/(3n)$. 这一结论根据三角不等式容易得到. 注意, 任意 $V\in\mathcal{A}$, 集合 $E_n(V)$ 包含于 $V$.


现在, 我们来定义
\begin{align*}
\mathcal{E}_n=\{E_n(U)\mid U\in\mathcal{A}\}.
\end{align*}
我们断言: $\mathcal{E}_n$ 是一个局部有限的开集族, 并且加细 $\mathcal{A}$. $\mathcal{E}_n$ 加细 $\mathcal{A}$ 是因为对于每一个 $V\in\mathcal{A}$, 有 $E_n(V)\subset V$. 再由任意 $x\in X$, $B(x,1/(6n))$ 最多与 $\mathcal{E}_n$ 中的一个元素相交, 可见 $\mathcal{E}_n$ 是局部有限的.

当然, $\mathcal{E}_n$ 并不是 $X$ 的一个覆盖. (\cref{figure:ch06-fig39-2}说明了这一点.) 然而我们可以断言: 集族
\begin{align*}
\mathcal{E}=\bigcup_{n\in\mathbb{Z}_+}\mathcal{E}_n
\end{align*}
覆盖 $X$.

设 $x\in X$, 族 $\mathcal{A}$ 如前所述是 $X$ 的一个覆盖. 选取 $U$ 为 $\mathcal{A}$ 中包含 $x$ 的第一个集合(按良序 $<$). 由于 $U$ 是一个开集, 所以可以选取适当的 $n$ 使得 $B(x,1/n)\subset U$. 于是, 根据定义易见 $x\in S_n(U)$. 因为 $U$ 是 $\mathcal{A}$ 中包含 $x$ 的第一个元素, 所以 $x\in T_n(U)$, 从而 $x$ 也属于 $\mathcal{E}_n$ 中的元素 $E_n(U)$.
\end{proof}

\subsection*{习题}

\begin{enumerate}
\item 验证例 1 中的结论.
\item 找出 $\mathbb{R}$ 的一个开覆盖 $\mathcal{A}$, 使得 $\mathcal{A}$ 是点态有限的, 却不是局部有限的. (集族 $\mathcal{A}$ 称为点态有限的, 若 $\mathbb{R}$ 中的每一个点只属于 $\mathcal{A}$ 中的有限多个元素.)
\item 试给出一集族 $\mathcal{A}$, 使得 $\mathcal{A}$ 本身不是局部有限的, 但 $\mathcal{B}=\{\overline{A}\mid A\in\mathcal{A}\}$ 是局部有限的.
\item 设 $\mathcal{A}$ 是 $\mathbb{R}$ 的子集族,
\begin{align*}
\mathcal{A}=\{(n,n+2)\mid n\in\mathbb{Z}\}.
\end{align*}
以下哪个子集族加细 $\mathcal{A}$?
\begin{align*}
\mathcal{B}&=\{(x,x+1)\mid x\in\mathbb{R}\},\\
\mathcal{C}&=\Bigl\{\bigl(n,n+\tfrac32\bigr)\mid n\in\mathbb{Z}\Bigr\},\\
\mathcal{D}&=\Bigl\{\bigl(x,x+\tfrac32\bigr)\mid x\in\mathbb{R}\Bigr\}.
\end{align*}
\item 证明: 如果 $X$ 有可数基, 那么 $X$ 的一个子集族 $\mathcal{A}$ 是可数局部有限的当且仅当它是可数的.
\item 在一致拓扑下考察 $\mathbb{R}^\omega$. 给定 $n$, 令 $\mathcal{B}_n$ 为 $\mathbb{R}^\omega$ 中所有形如 $\prod A_i$ 的子集族, 其中, 对于 $i\leqslant n$, $A_i=\mathbb{R}$; 对于其他情形, $A_i$ 等于 $\{0\}$ 或者 $\{1\}$. 证明: 集族 $\mathcal{B}=\bigcup\mathcal{B}_n$ 是可数局部有限的, 却既不是可数的, 也不是局部有限的.
\end{enumerate}

\end{document}